\documentclass[11pt]{article}

\usepackage[T1]{fontenc}
\usepackage[utf8]{inputenc}
\usepackage{lmodern}
\usepackage[margin=1in]{geometry}
\usepackage{amsmath,amssymb,amsthm}
\usepackage{booktabs}
\usepackage{array}
\usepackage{color}
\usepackage{hyperref}

\definecolor{linkblue}{rgb}{0,0,0.55}
\hypersetup{
  colorlinks=true,
  linkcolor=linkblue,
  citecolor=linkblue,
  urlcolor=linkblue
}

\newtheorem{theorem}{Theorem}[section]
\newtheorem{lemma}[theorem]{Lemma}
\newtheorem{proposition}[theorem]{Proposition}
\newtheorem{corollary}[theorem]{Corollary}
\theoremstyle{definition}
\newtheorem{definition}[theorem]{Definition}
\newtheorem{remark}[theorem]{Remark}
\newtheorem{example}[theorem]{Example}

\newcommand{\QQ}{\mathbb{Q}}
\newcommand{\RR}{\mathbb{R}}
\newcommand{\CC}{\mathbb{C}}
\newcommand{\ZZ}{\mathbb{Z}}
\newcommand{\NN}{\mathbb{N}}
\newcommand{\des}{\operatorname{des}}
\newcommand{\wdes}{\operatorname{wdes}}
\newcommand{\asc}{\operatorname{asc}}
\newcommand{\Li}{\operatorname{Li}}
\newcommand{\Bern}{\operatorname{Bernoulli}}
\newcommand{\Geom}{\operatorname{Geom}}

\title{\textbf{A Complete Disproof of Conjecture 00000000472:\\
Universal Real-Simplicity of Multi-Indexed Eulerian Polynomials}}
\author{Formalized Submission for Conjecture 00000000472}
\date{November 24, 2026}

\begin{document}
\maketitle

\begin{abstract}
Conjecture~00000000472 introduces $A_n(t)$ as the Eulerian polynomial of the Bernoulli
distribution and asserts that the multi-indexed Eulerian polynomial $A_\alpha(t)$
has all real simple roots if and only if the multi-index $\alpha$ is piecewise-constant.
We disprove this conjecture across every standard convention. Under the primary
combinatorial and generating-function definition---namely, the multiset Eulerian
polynomial $A_\alpha^{\mathrm{ms}}(t) = \sum_w t^{\des(w)}$ over multiset permutations
$w \in \mathfrak{S}(1^{\alpha_1}\cdots m^{\alpha_m})$ of MacMahon and Simion---we give an
elementary interlacing proof
that \emph{every} nonempty positive integer multi-index
$\alpha = (\alpha_1,\dots,\alpha_m) \in \ZZ_{\ge 1}^m$ yields a polynomial of exact degree
$|\alpha| - \max_i \alpha_i$ whose complex roots are all real, strictly negative, and
simple (and likewise for the shifted polynomial $\tilde{A}_\alpha^{\mathrm{ms}}(t) =
t A_\alpha^{\mathrm{ms}}(t)$). The same master differential operator proves
that the multiple polylogarithm Eulerian polynomials $A_\alpha^{\mathrm{mpl}}(t)$ and
$A_\alpha^{\mathrm{mpl},\star}(t)$ also have $|\alpha|-1$ distinct negative simple real
roots for \emph{every} $\alpha \in \ZZ_{\ge 1}^m$. In particular, strictly increasing
multi-indices with all run lengths equal to $1$---such as $\alpha=(1,2)$
with $A_{(1,2)}^{\mathrm{ms}}(t)=1+2t$, $\alpha=(1,2,3)$ with
$A_{(1,2,3)}^{\mathrm{ms}}(t)=1+17t+33t^2+9t^3$ ($\Delta=82{,}800>0$), and
the infinite family $\alpha^{(m)}=(1,2,\dots,m)$---all have exclusively real simple roots,
refuting necessity for every nontrivial definition of piecewise-constant index. Moreover,
under alternative conventions retaining $t^k$ prefactors or using Savage--Schuster
inversion-sequence Eulerian polynomials $E_\alpha(t)$, \emph{both} implications fail. All
main theorems, concrete witnesses, and non-vacuity checks are formalized from first
principles in Core Lean~4.19.0 (including constructive rationals $\QQ$, Cauchy reals $\RR$
with the polynomial Intermediate Value Theorem, and $\CC=\RR[i]$) with zero external
dependencies.
\end{abstract}

\section{Official Problem Statement and Semantic Analysis}

The official bilingual text of Conjecture~00000000472 (reproduced byte-for-byte in
\texttt{SOURCE.md}) states:

\begin{quote}
\textbf{English.} Definition: The Eulerian polynomial $A_n(t)$ of the Bernoulli
distribution. Conjecture: The multivariate Eulerian polynomial $A_\alpha(t)$ has all real
simple roots if and only if $\alpha$ is a piecewise-constant index (an explicit criterion,
programmable for small indices).
\end{quote}

The accompanying Chinese text clarifies the terminology precisely: after introducing the
Eulerian polynomial $A_n(t)$ of the Bernoulli distribution, it states the conjecture for
the \emph{multi-index Eulerian polynomial} $A_\alpha(t)$ (literally \emph{duochong zhibiao
Eulerian duoxiangshi $A_\alpha(t)$}), asserting that all of its roots are real simple roots
if and only if $\alpha$ is a piecewise-constant index (\emph{fenduan changshu zhibiao}, an
explicit criterion programmable for small indices). Notice that in both languages
$A_\alpha(t)$ is explicitly a \emph{univariate} polynomial in the single indeterminate $t$,
indexed by a finite multi-index $\alpha = (\alpha_1, \dots, \alpha_m) \in \ZZ_{\ge 1}^m$.

\subsection{Classical \texorpdfstring{$A_n(t)$}{A\_n(t)} and the Bernoulli Distribution}

Let $\mathfrak{S}_n$ denote the symmetric group on $\{1,\dots,n\}$. For $n \ge 1$, the
classical Eulerian polynomial is
\begin{equation}\label{eq:classical-eulerian}
  A_n(t) := \sum_{\pi \in \mathfrak{S}_n} t^{\des(\pi)}, \qquad
  \des(\pi) := \#\{1 \le i < n : \pi(i) > \pi(i+1)\},
\end{equation}
with $A_0(t) := 1$, so $A_1(t)=1$, $A_2(t)=1+t$, $A_3(t)=1+4t+t^2$, and Worpitzky's
power-sum / negative-polylogarithm identity holds for $|t| < 1$:
\begin{equation}\label{eq:worpitzky}
  \sum_{k=0}^{\infty} (k+1)^n t^k = \frac{A_n(t)}{(1-t)^{n+1}}, \qquad
  \Li_{-n}(t) := \sum_{k=1}^{\infty} k^n t^k = \frac{t A_n(t)}{(1-t)^{n+1}} \quad (n \ge 1).
\end{equation}
There are two canonical ways in which $A_n(t)$ governs the $\Bern(p)$ distribution:
\begin{enumerate}
  \item \textbf{Cumulants of $\Bern(p)$:} Let $X \sim \Bern(p)$ ($0 < p < 1$) with cumulant
  generating function $K_X(u) = \log \mathbb{E}[e^{uX}] = \log(1 - p + p e^u)$. Setting the
  negative odds parameter $t := -p/(1-p) \in (-\infty,0)$ gives $1 - p + p e^u =
  (1 - t e^u)/(1 - t)$. Differentiating $n+1$ times at $u = 0$ yields, for every $n \ge 1$,
  \begin{equation}\label{eq:bernoulli-cumulant}
    \kappa_{n+1}(X) = \left.\frac{d^{n+1}}{du^{n+1}} \log(1 - t e^u)\right|_{u=0}
    = -\Li_{-n}(t) = -\frac{t A_n(t)}{(1-t)^{n+1}}
    = p(1-p)^n A_n\!\left(-\frac{p}{1-p}\right).
  \end{equation}
  Equivalently, by the chain rule $dp(u)/du = p(u)(1-p(u))$ for $p(u) = pe^u/(1-p+pe^u)$,
  the cumulant polynomials satisfy $\kappa_1(p) = p$ and $\kappa_{n+1}(p) = p(1-p)\kappa_n'(p)$.
  \item \textbf{Waiting-time moments of Bernoulli trials:} If $G \sim \Geom(1-p)$ counts the
  number of independent $\Bern(1-p)$ trials up to and including the first success, then
  $\mathbb{E}[G^n] = (1-p)\sum_{k=0}^\infty (k+1)^n p^k = A_n(p)/(1-p)^n$.
\end{enumerate}
By a classical theorem of Frobenius~\cite{Frobenius1910}, $A_n(t)$ has $\max(0,n-1)$
distinct negative simple real roots for every $n \ge 1$.

\subsection{Standard Multi-Index Generalizations \texorpdfstring{$A_\alpha(t)$}{A\_alpha(t)}}

Let $\alpha = (\alpha_1, \dots, \alpha_m) \in \ZZ_{\ge 1}^m$ ($m \ge 1$) be a positive
multi-index of total weight $|\alpha| := \alpha_1 + \dots + \alpha_m$ and maximum entry
$M(\alpha) := \max_{1 \le i \le m} \alpha_i$. Across algebraic combinatorics, analytic
number theory, and probability, there are four univariate multi-index extensions of $A_n(t)$:

\begin{definition}[Multiset Eulerian Polynomial~\cite{MacMahon1915,DillonRoselle1969,Simion1984}]\label{def:multiset-eulerian}
For $\alpha \in \ZZ_{\ge 1}^m$, let $\mathcal{M}(\alpha) = \{1^{\alpha_1}, \dots, m^{\alpha_m}\}$
be the multiset with $\alpha_i$ copies of $i$, and let $\mathfrak{S}(\mathcal{M}(\alpha))$ denote
the set of all $|\alpha|! / (\alpha_1! \cdots \alpha_m!)$ distinct words of length $|\alpha|$
with letter multiplicities $\alpha_1,\dots,\alpha_m$. The \emph{multiset Eulerian polynomial}
is
\begin{equation}\label{eq:ms-def}
  A_\alpha^{\mathrm{ms}}(t) := \sum_{w \in \mathfrak{S}(\mathcal{M}(\alpha))} t^{\des(w)},
  \qquad \des(w) := \#\{1 \le j < |\alpha| : w_j > w_{j+1}\},
\end{equation}
which satisfies MacMahon's rational generating function identity
\begin{equation}\label{eq:macmahon}
  \sum_{k=0}^{\infty} \left(\prod_{i=1}^{m} \binom{k + \alpha_i}{\alpha_i}\right) t^k
  = \frac{A_\alpha^{\mathrm{ms}}(t)}{(1-t)^{|\alpha|+1}}.
\end{equation}
When $\alpha = (1,\dots,1) \in \ZZ_{\ge 1}^n$, $A_{(1,\dots,1)}^{\mathrm{ms}}(t) = A_n(t)$.
One also considers the $t$-shifted convention $\tilde{A}_\alpha^{\mathrm{ms}}(t) :=
t A_\alpha^{\mathrm{ms}}(t) = \sum_w t^{\des(w)+1}$ and the weak-descent polynomial
$W_\alpha(t) := \sum_{w \in \mathfrak{S}(\mathcal{M}(\alpha))} t^{\wdes(w)}$ where
$\wdes(w) := \#\{1 \le j < |\alpha| : w_j \ge w_{j+1}\}$.
\end{definition}

\begin{definition}[Multiple Polylogarithm Eulerian Polynomials~\cite{TaNgo2026}]\label{def:mpl-eulerian}
Generalizing~\eqref{eq:worpitzky} to depth-$m$ multiple polylogarithms at negative integer
indices, for $|t| < 1$ we have rational generating functions
\begin{align}
  \Li_{-\alpha_1,\dots,-\alpha_m}(t)
  &:= \sum_{n_1 > n_2 > \dots > n_m \ge 1} n_1^{\alpha_1}\cdots n_m^{\alpha_m} t^{n_1}
  = \frac{t^m A_\alpha^{\mathrm{mpl}}(t)}{(1-t)^{|\alpha|+m}}, \label{eq:mpl-strict}\\
  \Li_{-\alpha_1,\dots,-\alpha_m}^{\star}(t)
  &:= \sum_{n_1 \ge n_2 \ge \dots \ge n_m \ge 1} n_1^{\alpha_1}\cdots n_m^{\alpha_m} t^{n_1}
  = \frac{t A_\alpha^{\mathrm{mpl},\star}(t)}{(1-t)^{|\alpha|+m}}, \label{eq:mpl-star}
\end{align}
where $A_\alpha^{\mathrm{mpl}}(t)$ and $A_\alpha^{\mathrm{mpl},\star}(t)$ have nonzero
constant terms and recover $A_{(n)}^{\mathrm{mpl}}(t) = A_{(n)}^{\mathrm{mpl},\star}(t) =
A_n(t)$ at depth $m=1$.
\end{definition}

\begin{definition}[Other Multi-Index Conventions]\label{def:other-conventions}
Two further candidates are: (i)~the \emph{comonotonic joint Bernoulli cumulant}
$\kappa_\alpha(X,\dots,X) = \kappa_{|\alpha|}(X)$ for $X \sim \Bern(p)$, whose numerator is
$A_{|\alpha|-1}(t) = A_{(1,\dots,1)}^{\mathrm{ms}}(t)$; and (ii)~the \emph{Savage--Schuster
$\alpha$-Eulerian polynomial}~\cite{SavageSchuster2012,SavageVisontai2015}
$E_\alpha(t) := \sum_{e \in \mathcal{I}_m(\alpha)} t^{\asc_\alpha(e)}$ over inversion
sequences $\mathcal{I}_m(\alpha) = \prod_{i=1}^m \{0,\dots,\alpha_i-1\}$ with
$\asc_\alpha(e) := \#\{0 \le i < m : e_i/\alpha_i < e_{i+1}/\alpha_{i+1}\}$ ($e_0=0,\alpha_0=1$).
\end{definition}

\begin{definition}[AllRealSimpleRoots and Piecewise-Constant Multi-Indices]\label{def:predicates}
A nonzero polynomial $P(t) \in \RR[t]$ \emph{has all real simple roots} if every complex
root $z \in \CC$ of $P(z)=0$ satisfies $\operatorname{Im}(z)=0$ and $P'(z)\neq 0$
(equivalently, $P(t) = c \prod_{j=1}^d (t - r_j)$ for some $c \in \RR \setminus \{0\}$ and
pairwise distinct $r_1,\dots,r_d \in \RR$; note that a nonzero constant polynomial of degree
$d=0$ vacuously has all complex roots real and simple).

For a finite multi-index $\alpha = (\alpha_1,\dots,\alpha_m) \in \ZZ_{\ge 1}^m$, standard
nontrivial interpretations of ``$\alpha$ is a piecewise-constant index'' include:
\begin{enumerate}
  \item \textbf{Constant index:} $\alpha_1 = \alpha_2 = \dots = \alpha_m$.
  \item \textbf{Block-constant index (minimum run length $\ge 2$):} $\alpha$ decomposes into
  constant contiguous runs each of length at least $2$ (e.g.\ $(1,1,2,2)$ or $(2,2,2)$).
  \item \textbf{Has at least one constant run ($\exists i,\ \alpha_i = \alpha_{i+1}$):}
  not all adjacent entries are distinct.
  \item \textbf{Arbitrary proper criterion:} \emph{any} predicate $\mathcal{P} \subsetneq
  \ZZ_{\ge 1}^+$ that excludes at least one valid multi-index (for a vacuous predicate
  $\mathcal{P} = \ZZ_{\ge 1}^+$ is not a discriminating criterion).
\end{enumerate}
Crucially, every strictly increasing multi-index $\alpha^{(m)} := (1,2,\dots,m)$ ($m \ge 2$)
---in particular $\alpha=(1,2)$ and $\alpha=(1,2,3)$---has pairwise distinct entries, all
constant runs of length $1$, and zero adjacent equalities, failing (1)--(3) simultaneously.
\end{definition}

\section{Combinatorial Insertion Reciprocity and Simion's Recurrence}

We establish Simion's differential recurrence~\cite{Simion1984} directly from the
combinatorial definition~\eqref{eq:ms-def} via a finite insertion--deletion bijection.

\begin{lemma}[Insertion--Deletion Reciprocity and Descent Transition]\label{lem:insertion}
Let $m \ge 1$, and let $w = (w_1,\dots,w_N) \in \{1,\dots,m\}^N$ be any word whose entries
satisfy $w_j \le m$ for all $1 \le j \le N$, containing $c := \#\{j : w_j = m\}$ copies of
the maximum letter $m$ and $k := \des(w)$ strict descents. For $0 \le p \le N$, let
$\operatorname{ins}_p(m, w)$ denote the word of length $N+1$ obtained by inserting $m$ after
the first $p$ entries of $w$. Then:
\begin{enumerate}
  \item Among the $N+1$ words $\{\operatorname{ins}_p(m,w)\}_{p=0}^N$, exactly $c + 1 + k$
  have $\des(\operatorname{ins}_p(m,w)) = k$, and the remaining $N - c - k$ have
  $\des(\operatorname{ins}_p(m,w)) = k + 1$. Consequently,
  \begin{equation}\label{eq:word-sum}
    \sum_{p=0}^{N} t^{\des(\operatorname{ins}_p(m,w))}
    = (c + 1 + k)\,t^k + (N - c - k)\,t^{k+1}.
  \end{equation}
  \item For every word $w' \in \{1,\dots,m\}^{N+1}$, the number of indices $p \in \{0,\dots,N\}$
  with $\operatorname{ins}_p(m,w) = w'$ equals the number of indices $p \in \{0,\dots,N\}$
  with $w'_{p+1} = m$ whose deletion from $w'$ yields $w$.
\end{enumerate}
\end{lemma}

\begin{proof}
For (1), inspect each insertion gap $p \in \{0,\dots,N\}$:
\begin{itemize}
  \item At the right endpoint $p = N$, $w_N \le m$ creates no descent, so
  $\des(\operatorname{ins}_N(m,w)) = k$ ($1$ position).
  \item At an interior gap $1 \le p \le N-1$: the pair $(w_p, w_{p+1})$ is replaced by
  $(w_p, m)$ and $(m, w_{p+1})$. Since $w_p \le m$, $(w_p, m)$ is never a strict descent,
  while $(m, w_{p+1})$ is a strict descent if and only if $w_{p+1} < m$. Hence:
  if $w_p > w_{p+1}$ (so automatically $w_{p+1} < w_p \le m$), one old descent is lost and
  one new descent is gained, giving $\des = k$; if $w_p \le w_{p+1}$ and $w_{p+1} = m$, no
  descent is lost or gained, giving $\des = k$; if $w_p \le w_{p+1} < m$, one new descent is
  gained, giving $\des = k+1$.
  \item At the left endpoint $p = 0$, $\des(\operatorname{ins}_0(m,w)) = k$ if $w_1 = m$,
  and $k+1$ if $w_1 < m$.
\end{itemize}
Because $w_p > m$ is impossible, no descent $w_p > w_{p+1}$ can satisfy $w_{p+1} = m$. Thus
the $k$ interior descent gaps $\{1 \le p < N : w_p > w_{p+1}\}$ and the $c$ gaps
$\{0 \le p < N : w_{p+1} = m\}$ immediately preceding a copy of $m$ are disjoint, and
neither contains $p=N$. Hence $1 + k + c$ gaps yield $\des = k$ and $(N+1) - (c+1+k) =
N - c - k$ gaps yield $\des = k+1$, proving~\eqref{eq:word-sum}. Part~(2) is immediate
because $\operatorname{ins}_p(m,w) = w'$ holds if and only if $w'_{p+1} = m$ and deleting
index $p+1$ from $w'$ recovers $w$.
\end{proof}

\begin{proposition}[Simion's Operator Recurrence]\label{prop:simion-recurrence}
Define the linear differential operator $\mathcal{T}_{r,s} : \RR[t] \to \RR[t]$ by
\begin{equation}\label{eq:interlace-op}
  \mathcal{T}_{r,s}(P)(t) := t(1-t)P'(t) + (r + st)P(t).
\end{equation}
Let $\beta = (\alpha_1,\dots,\alpha_{m-1}, r-1)$ ($r \ge 1$, where $r-1=0$ denotes
$(\alpha_1,\dots,\alpha_{m-1})$) have total weight $N = |\beta|$, and let
$\alpha = (\alpha_1,\dots,\alpha_{m-1}, r)$. Then
\begin{equation}\label{eq:simion-rec}
  r\,A_\alpha^{\mathrm{ms}}(t) = \mathcal{T}_{r,\,N+1-r}\!\bigl(A_\beta^{\mathrm{ms}}\bigr)(t).
\end{equation}
\end{proposition}

\begin{proof}
Every $w' \in \mathfrak{S}(\mathcal{M}(\alpha))$ contains exactly $r$ copies of $m$, while
every $w \in \mathfrak{S}(\mathcal{M}(\beta))$ has length $N$, entries $\le m$, and
$c = r - 1$ copies of $m$. Summing $t^{\des(\operatorname{ins}_p(m,w))}$ over all
$w \in \mathfrak{S}(\mathcal{M}(\beta))$ and $0 \le p \le N$, Lemma~\ref{lem:insertion}(2)
counts each $w' \in \mathfrak{S}(\mathcal{M}(\alpha))$ exactly $r$ times, whereas
Lemma~\ref{lem:insertion}(1) evaluates the inner sum as $(r + \des(w))t^{\des(w)} +
(N + 1 - r - \des(w))t^{\des(w)+1}$. Summing over $w$ gives
\[
  r\,A_\alpha^{\mathrm{ms}}(t)
  = r A_\beta^{\mathrm{ms}}(t) + t A_\beta^{\mathrm{ms}\prime}(t)
    + (N+1-r)t A_\beta^{\mathrm{ms}}(t) - t^2 A_\beta^{\mathrm{ms}\prime}(t)
  = \mathcal{T}_{r,\,N+1-r}\!\bigl(A_\beta^{\mathrm{ms}}\bigr)(t).\qedhere
\]
\end{proof}

\section{Master Interlacing Theorem and Universal Real-Simplicity}

We now prove our central root-simplicity theorem by a purely algebraic sign-alternation
argument coupled with the polynomial Intermediate Value Theorem.

\begin{theorem}[Master Interlacing Lemma]\label{thm:master-interlacing}
Suppose $P(t) = c \prod_{j=1}^d (t - \xi_j) \in \RR[t]$ with $c > 0$ and $d \ge 0$ strictly
decreasing negative roots $0 > \xi_1 > \xi_2 > \dots > \xi_d$. Let $r > 0$ and $s \ge d$ be
real numbers, and set $Q := \mathcal{T}_{r,s}(P)$.
\begin{enumerate}
  \item \textbf{Degree-preserving case ($s = d$):} $Q$ has exact degree $d$ and splits as
  $Q(t) = c' \prod_{j=1}^d (t - \eta_j)$ with $c' > 0$ and $d$ strictly decreasing negative
  roots satisfying $0 > \eta_1 > \xi_1 > \eta_2 > \xi_2 > \dots > \eta_d > \xi_d$.
  \item \textbf{Degree-increasing case ($s > d$):} $Q$ has exact degree $d+1$ with leading
  coefficient $c' = (s - d)c > 0$, and splits as $Q(t) = c' \prod_{j=1}^{d+1} (t - \eta_j)$
  with $d+1$ strictly decreasing negative roots satisfying
  $0 > \eta_1 > \xi_1 > \eta_2 > \dots > \eta_d > \xi_d > \eta_{d+1}$.
\end{enumerate}
\end{theorem}

\begin{proof}
Write $P(t) = c t^d + R(t)$ with $\deg R \le d-1$. In
$Q(t) = t(1-t)P'(t) + (r+st)P(t)$, the coefficient of $t^{d+1}$ is $(-d + s)c = (s-d)c$, and
all powers $t^k$ with $k > d+1$ vanish. Thus $\deg Q \le d$ when $s=d$, whereas $\deg Q = d+1$
with positive leading coefficient $c' = (s-d)c > 0$ when $s > d$.

Set $\xi_0 := 0$. At $t = 0$, $P(0) = c\prod_{j=1}^d (-\xi_j) > 0$ since $-\xi_j > 0$, so
$Q(\xi_0) = r P(0) > 0$. For $d=0$ and $s=0$, $Q(t) = rc > 0$ immediately satisfies~(1).
Otherwise, for each $k \in \{1,\dots,d\}$, $P(\xi_k) = 0$ and the formal product rule gives
\[
  P'(\xi_k) = c \prod_{j \neq k} (\xi_k - \xi_j).
\]
Among $j \neq k$, the $k-1$ indices $j < k$ satisfy $\xi_k - \xi_j < 0$ and the $d-k$ indices
$j > k$ satisfy $\xi_k - \xi_j > 0$, so $(-1)^{k-1} P'(\xi_k) > 0$. Since $\xi_k < 0$ implies
$\xi_k(1-\xi_k) < 0$, we obtain
\begin{equation}\label{eq:alt-signs}
  (-1)^k Q(\xi_k) = (-1)^k \xi_k(1-\xi_k) P'(\xi_k) > 0
  \qquad \text{for every } k \in \{0, 1, \dots, d\}.
\end{equation}
Thus $Q$ strictly alternates in sign across the $d+1$ points
$\xi_d < \xi_{d-1} < \dots < \xi_1 < \xi_0 = 0$. By the Intermediate Value Theorem on each
open subinterval $(\xi_k, \xi_{k-1})$ ($1 \le k \le d$), there exist $d$ roots $\eta_k \in
(\xi_k, \xi_{k-1})$ of $Q$, satisfying $0 > \eta_1 > \xi_1 > \dots > \eta_d > \xi_d$.

In case~(1) ($s = d$, $d \ge 1$), $\deg Q \le d$ and $Q$ vanishes at $d$ distinct points, so
$Q(t) = c' \prod_{j=1}^d (t - \eta_j)$ for some $c' \in \RR$. Evaluating at $t = 0$ gives
$0 < Q(0) = c' \prod_{j=1}^d (-\eta_j)$, forcing $c' > 0$ since each $-\eta_j > 0$.

In case~(2) ($s > d$), $Q(t) = c' t^{d+1} + \sum_{i=0}^d q_i t^i$ with $c' = (s-d)c > 0$.
Choosing $L < \min(-1, \xi_d, -2(|q_0|+\dots+|q_d|)/c')$, we have
$|q_d L^{-1} + \dots + q_0 L^{-(d+1)}| \le (|q_0|+\dots+|q_d|)/|L| < c'/2$, whence
$(-1)^{d+1} Q(L) = |L|^{d+1}(c' + q_d L^{-1} + \dots + q_0 L^{-(d+1)}) > |L|^{d+1}(c'/2) > 0$.
Comparing with $(-1)^d Q(\xi_d) > 0$ from~\eqref{eq:alt-signs}, the Intermediate Value
Theorem on $(L, \xi_d)$ yields an additional root $\eta_{d+1} \in (L, \xi_d) \subset
(-\infty, \xi_d)$. Hence $Q$ has $d+1 = \deg Q$ distinct negative roots
$0 > \eta_1 > \dots > \eta_{d+1}$ and factors as $Q(t) = c'\prod_{j=1}^{d+1}(t-\eta_j)$.
\end{proof}

\begin{lemma}[Factorization Criterion for Simple Real Roots]\label{lem:simple-roots}
If $P(t) = c \prod_{j=1}^d (t - r_j) \in \RR[t]$ with $c \neq 0$ and pairwise distinct real
numbers $r_1,\dots,r_d \in \RR$, then every root $z \in \CC$ of $P(z)=0$ satisfies
$\operatorname{Im}(z) = 0$ and $P'(z) \neq 0$.
\end{lemma}

\begin{proof}
Over the field $\CC = \RR[i]$, $P(z) = c \prod_{j=1}^d (z - r_j) = 0$ implies $z = r_k \in \RR$
for some $k$, so $\operatorname{Im}(z) = 0$. Moreover, $P'(r_k) = c \prod_{j \neq k}
(r_k - r_j) \neq 0$ because $c \neq 0$ and $r_k \neq r_j$ for all $j \neq k$.
\end{proof}

\begin{theorem}[Universal Real-Simplicity of Multiset Eulerian Polynomials]\label{thm:universal-ms}
For \emph{every} nonempty positive integer multi-index
$\alpha = (\alpha_1,\dots,\alpha_m) \in \ZZ_{\ge 1}^m$:
\begin{enumerate}
  \item $A_\alpha^{\mathrm{ms}}(t)$ has exact degree $d(\alpha) = |\alpha| - \max_{1\le i\le m}
  \alpha_i$, positive leading coefficient, constant term $A_\alpha^{\mathrm{ms}}(0)=1$, and
  splits over $\RR$ into $d(\alpha)$ distinct linear factors with strictly negative roots.
  In particular, $A_\alpha^{\mathrm{ms}}(t)$ has all real simple roots for \emph{every}
  $\alpha \in \ZZ_{\ge 1}^m$.
  \item The shifted polynomial $\tilde{A}_\alpha^{\mathrm{ms}}(t) = t A_\alpha^{\mathrm{ms}}(t)$
  has exact degree $d(\alpha)+1$ and $d(\alpha)+1$ distinct non-positive real simple roots.
\end{enumerate}
\end{theorem}

\begin{proof}
We induct on the construction of $\alpha$. For $m=1$, $\mathfrak{S}(1^{\alpha_1}) =
\{(1,\dots,1)\}$ has $\des = 0$, so $A_{(\alpha_1)}^{\mathrm{ms}}(t) = 1$, which has degree
$0 = \alpha_1 - \alpha_1$ and $0$ roots. Now suppose the claim holds for
$\beta = (\alpha_1,\dots,\alpha_{m-1}, r-1)$ ($r \ge 1$) of weight $N = |\beta|$ and maximum
$M = \max(\alpha_1,\dots,\alpha_{m-1}, r-1)$, so $\deg A_\beta^{\mathrm{ms}} = d = N - M$.
By Proposition~\ref{prop:simion-recurrence}, $A_\alpha^{\mathrm{ms}} =
\frac{1}{r}\mathcal{T}_{r,s}(A_\beta^{\mathrm{ms}})$ with $s = N + 1 - r$. Note that
$s - d = (N + 1 - r) - (N - M) = M + 1 - r$:
\begin{itemize}
  \item If $r \le M$, then $\max(\alpha) = M$, $s \ge d + 1 > d$, and
  Theorem~\ref{thm:master-interlacing}(2) shows that $A_\alpha^{\mathrm{ms}}$ has degree
  $d + 1 = (N+1) - M$ and splits into $d+1$ distinct negative real linear factors.
  \item If $r = M + 1$ (which occurs only when $r - 1 = M$), then $\max(\alpha) = M + 1$,
  $s = d$, and Theorem~\ref{thm:master-interlacing}(1) shows that $A_\alpha^{\mathrm{ms}}$
  has degree $d = (N+1) - (M+1)$ and splits into $d$ distinct negative real linear factors.
\end{itemize}
Lemma~\ref{lem:simple-roots} yields real simplicity of all complex roots. Finally, since all
roots of $A_\alpha^{\mathrm{ms}}$ are strictly negative, $t=0$ is not a root of
$A_\alpha^{\mathrm{ms}}$, so $t A_\alpha^{\mathrm{ms}}(t)$ also has pairwise distinct real
roots, proving~(2).
\end{proof}

\begin{theorem}[Universal Real-Simplicity of Multiple Polylogarithm Eulerian Polynomials]\label{thm:universal-mpl}
For \emph{every} $\alpha = (\alpha_1,\dots,\alpha_m) \in \ZZ_{\ge 1}^m$, both
$A_\alpha^{\mathrm{mpl}}(t)$ and $A_\alpha^{\mathrm{mpl},\star}(t)$ have exact degree
$|\alpha| - 1$ and split over $\RR$ into $|\alpha| - 1$ distinct linear factors with strictly
negative real roots.
\end{theorem}

\begin{proof}
From~\eqref{eq:mpl-strict}, $\Li_{-\alpha_1,\dots,-\alpha_m}(t) = (t\frac{d}{dt})^{\alpha_1}
\bigl(\frac{t}{1-t}\Li_{-\alpha_2,\dots,-\alpha_m}(t)\bigr)$. Since
\[
  t\frac{d}{dt}\left(\frac{t^k P(t)}{(1-t)^M}\right)
  = \frac{t^k\,\mathcal{T}_{k,\,M-k}(P)(t)}{(1-t)^{M+1}},
\]
starting at the base $\alpha_m = 1$ where $\Li_{-1}(t) = \frac{t \cdot 1}{(1-t)^2}$ ($k=1,
M=2, \deg P = 0$), every subsequent application of $t\frac{d}{dt}$ (incrementing $M$ while
fixing $k$) or prepending of $\frac{t}{1-t}$ followed by $t\frac{d}{dt}$ (incrementing both
$k$ and $M$, leaving $M-k$ equal to the cumulative weight $|\text{suffix}| = \deg P + 1 >
\deg P$) invokes Theorem~\ref{thm:master-interlacing}(2) with $s = \deg P + 1 > \deg P$.
Similarly,~\eqref{eq:mpl-star} keeps $k=1$ throughout and applies $\mathcal{T}_{1,s}$ with
$s = M - 1 > \deg P$ at every step after the initial step. By induction, both polynomials
have degree $|\alpha|-1$ and $|\alpha|-1$ distinct negative real roots.
\end{proof}

\section{Explicit Counterexamples and Bidirectional Refutation}

\subsection{Explicit Non-Piecewise-Constant Counterexamples}

\begin{proposition}[Concrete Counterexamples]\label{prop:concrete}
\begin{enumerate}
  \item For $\alpha = (1,2)$, the $3!/(1!2!)=3$ permutations of $\{1,2,2\}$ are
  $122$ ($\des=0$), $212$ ($\des=1$), and $221$ ($\des=1$), so
  $A_{(1,2)}^{\mathrm{ms}}(t) = 1 + 2t$, which has the unique simple real root $t = -1/2$.
  \item For $\alpha = (1,2,3)$, the $\frac{6!}{1!2!3!} = 60$ multiset permutations of
  $\{1,2,2,3,3,3\}$ yield
  \begin{equation}\label{eq:A123}
    A_{(1,2,3)}^{\mathrm{ms}}(t) = 1 + 17t + 33t^2 + 9t^3.
  \end{equation}
  Evaluating at four rational points $-10 < -1 < -\frac{1}{3} < 0$ produces strict sign
  alternation:
  \begin{align*}
    A_{(1,2,3)}^{\mathrm{ms}}(-10) &= -5869 < 0, &
    A_{(1,2,3)}^{\mathrm{ms}}(-1) &= 8 > 0, \\
    A_{(1,2,3)}^{\mathrm{ms}}\!\left(-\tfrac{1}{3}\right) &= -\tfrac{4}{3} < 0, &
    A_{(1,2,3)}^{\mathrm{ms}}(0) &= 1 > 0,
  \end{align*}
  isolating three distinct simple real roots in $(-10,-1)$, $(-1,-1/3)$, and $(-1/3,0)$; its
  cubic discriminant is $\Delta = 33^2\cdot 17^2 - 4\cdot 9\cdot 17^3 - 4\cdot 33^3\cdot 1
  - 27\cdot 9^2\cdot 1^2 + 18\cdot 9\cdot 33\cdot 17\cdot 1 = 82{,}800 > 0$.
  \item For every $m \ge 2$, the strictly increasing index $\alpha^{(m)} = (1,2,\dots,m)$ has
  $m$ pairwise distinct entries (maximum constant-run length $1$), while
  $A_{\alpha^{(m)}}^{\mathrm{ms}}(t)$ has degree $\binom{m}{2} \ge 1$ and $\binom{m}{2}$
  distinct negative simple real roots.
\end{enumerate}
\end{proposition}

\subsection{Complete Refutation Across All Conventions}

\begin{corollary}[Disproof of Conjecture~00000000472]\label{cor:disproof}
Conjecture~00000000472 is false under every interpretation:
\begin{enumerate}
  \item \textbf{Standard conventions ($A_\alpha^{\mathrm{ms}}$, $\tilde{A}_\alpha^{\mathrm{ms}}$,
  $A_\alpha^{\mathrm{mpl}}$, $A_\alpha^{\mathrm{mpl},\star}$, and $\kappa_{|\alpha|+1}$):}
  By Theorems~\ref{thm:universal-ms} and~\ref{thm:universal-mpl}, \emph{every} valid
  multi-index $\alpha \in \ZZ_{\ge 1}^m$ produces all real simple roots. Hence for \emph{every}
  proper subset $\mathcal{P} \subsetneq \ZZ_{\ge 1}^+$ (in particular, any definition of
  piecewise-constant index excluding $\alpha=(1,2)$ or $\alpha=(1,2,3)$), the biconditional
  \[
    \forall \alpha \in \ZZ_{\ge 1}^+,\quad
    \operatorname{AllRealSimpleRoots}(A_\alpha) \Longleftrightarrow \alpha \in \mathcal{P}
  \]
  fails in the ``only if'' direction.
  \item \textbf{Alternative conventions (bidirectional failure):}
  \begin{itemize}
    \item \emph{Unstripped weak-descent polynomial} $W_\alpha(t) = t^{|\alpha|-1}
    A_\alpha^{\mathrm{ms}}(t^{-1})$: non-piecewise-constant $\alpha=(1,2)$ gives
    $W_{(1,2)}(t) = 2t + t^2 = t(2+t)$ with simple real roots $0,-2$, whereas constant
    $\alpha=(3,3)$ gives $W_{(3,3)}(t) = t^2(1 + 9t + 9t^2 + t^3)$ with a double root at $t=0$.
    \item \emph{Unstripped polylogarithm numerator} $P_\alpha^{\mathrm{mpl}}(t) =
    t^m A_\alpha^{\mathrm{mpl}}(t)$: non-constant $\alpha=(2)$ gives $t(1+t)$ (simple roots),
    whereas constant $\alpha=(1,1)$ gives $t^2(2+t)$ (double root at $t=0$).
    \item \emph{Savage--Schuster $\alpha$-Eulerian polynomial} $E_\alpha(t)$: strictly
    increasing $\alpha=(1,2,3)$ gives $E_{(1,2,3)}(t) = 1 + 4t + t^2$ (two simple real roots),
    whereas block-constant $\alpha=(2,2,1,1,2,2)$ (all runs of length $2$) gives
    $E_{(2,2,1,1,2,2)}(t) = (1 + 3t)^2$, which has a double root at $t = -1/3$.
  \end{itemize}
\end{enumerate}
\end{corollary}

\section{Formal Verification in Core Lean 4.19.0 and Computational Audit}

The entire proof is formalized from first principles in \texttt{lean/Main.lean} within
the namespace \texttt{Conjecture00000000472}, using Core Lean~\texttt{v4.19.0} with
\texttt{-DwarningAsError=true} and zero external dependencies:
\begin{itemize}
  \item \textbf{Constructive Field Tower $\QQ \subset \RR \subset \CC$:} \texttt{Q} is built
  as the quotient of $\ZZ \times \ZZ_{>0}$; \texttt{R} is built as the Cauchy completion of
  \texttt{Q} with dyadic moduli $2^{-k}$, proved to be an Archimedean ordered field in which
  $\QQ$ is dense (\texttt{R.dense\_ofQ}); and \texttt{Complex} is constructed as $\RR \times \RR$
  and proved to be an integral domain (\texttt{Complex.mul\_eq\_zero}).
  \item \textbf{Polynomial IVT and Root-Simplicity Equivalence:} Polynomials over $\RR$,
  Horner derivatives, synthetic linear division, Lipschitz bounds on compact intervals, and
  bisection over $\QQ$ yield \texttt{Poly.poly\_ivt}, root factorization
  (\texttt{Poly.factor\_of\_nodup\_roots}), and complex root-simplicity criteria.
  \item \textbf{Eulerian Families:} Duplicate-free multiset words (\texttt{multisetPerms},
  certified exact via letter-multiplicity histograms), \texttt{eulerianPoly}
  ($A_\alpha^{\mathrm{ms}}$), the insertion--deletion reciprocity proof of
  \texttt{simion\_recurrence}, and \texttt{mplEulerianPoly} ($A_\alpha^{\mathrm{mpl}}$).
  \item \textbf{Main Theorems and Axiom Audits:} Universal real-simplicity theorems for
  multi-indexed Eulerian families, explicit evaluations at $\alpha=(1,2)$ and
  $\alpha=(1,2,3)$, bidirectional counterexamples under alternative conventions, non-vacuity
  witnesses, and axiom checks confirming reliance solely on \texttt{propext},
  \texttt{Classical.choice}, and \texttt{Quot.sound}.
\end{itemize}

\begin{table}[htbp]
\centering
\footnotesize
\renewcommand{\arraystretch}{1.15}
\begin{tabular}{ll}
\hline
\textbf{Mathematical Result} & \textbf{Lean 4 Declaration (\texttt{Main.lean})} \\
\hline
Polynomial Intermediate Value Theorem & \texttt{Poly.poly\_ivt} \\
Master Interlacing Theorem (Thm.~\ref{thm:master-interlacing}) & \texttt{Poly.interlaceOp\_preserves\_negSimpleRoots} \\
Simion's Recurrence (Prop.~\ref{prop:simion-recurrence}) & \texttt{simion\_recurrence} \\
Universal Real-Simplicity of $A_\alpha^{\mathrm{ms}}, \tilde{A}_\alpha^{\mathrm{ms}}$ (Thm.~\ref{thm:universal-ms}) & \texttt{multiset\_allRealSimple} \\
Universal Real-Simplicity of $A_\alpha^{\mathrm{mpl}}$ (Thm.~\ref{thm:universal-mpl}) & \texttt{mpl\_allRealSimple} \\
Universal Real-Simplicity of $A_\alpha^{\mathrm{mpl},\star}$ (Thm.~\ref{thm:universal-mpl}) & \texttt{mplStar\_allRealSimple} \\
$\alpha=(1,2,3)$ Certificate \& Discriminant (Prop.~\ref{prop:concrete}) & \texttt{counterexample\_123\_allRealSimple} \\
Infinite Strictly Increasing Family $\alpha^{(m)}$ & \texttt{strictIncFamily\_counterexamples} \\
Refutation for Any Proper Criterion (Cor.~\ref{cor:disproof}) & \texttt{disproof\_conjecture472\_multiset\_anyProper} \\
Weak-Descent Bidirectional Failure & \texttt{weakDescent\_bidirectional\_failure} \\
Savage--Schuster Bidirectional Failure & \texttt{savageSchuster\_bidirectional\_failure} \\
Non-Vacuity Witness ($\alpha=(1,1)$) & \texttt{positive\_witness\_constant\_11} \\
Non-Tautology Witness ($1+t^2$ vs.\ $\CC$-roots) & \texttt{positive\_witness\_nonreal\_roots} \\
\hline
\end{tabular}
\caption{Correspondence between paper theorems and Core Lean 4.19.0 declarations.}
\label{tab:lean-map}
\end{table}

Independently, \texttt{verify.py} verifies via exact arbitrary-precision rational Sturm
chains that $511/511$ multiset Eulerian polynomials ($|\alpha|\le 9$) and $255/255$ multiple
polylogarithm Eulerian polynomials ($|\alpha|\le 8$) have exclusively simple negative real
roots, and checks the $3$-way identity between multiset permutations, MacMahon's rational
series, and Simion's recurrence across all $127$ multi-indices with $|\alpha|\le 7$.

\begin{thebibliography}{9}

\bibitem{DillonRoselle1969}
J.~F. Dillon and D.~P. Roselle,
\emph{Simon Newcomb's problem},
SIAM J.\ Appl.\ Math.\ \textbf{17} (1969), 1086--1093.

\bibitem{Frobenius1910}
G.~Frobenius,
\emph{\"Uber die Bernoullischen Zahlen und die Eulerschen Polynome},
Sitzungsber.\ Preuss.\ Akad.\ Wiss.\ Berlin (1910), 809--847.

\bibitem{MacMahon1915}
P.~A. MacMahon,
\emph{Combinatory Analysis}, Vols.~I--II,
Cambridge University Press, Cambridge, 1915--1916.

\bibitem{SavageSchuster2012}
C.~D. Savage and M.~J. Schuster,
\emph{Ehrhart series of lecture hall polytopes and Eulerian polynomials for inversion sequences},
J.\ Combin.\ Theory Ser.\ A \textbf{119} (2012), 850--870.

\bibitem{SavageVisontai2015}
C.~D. Savage and M.~Visontai,
\emph{The $\mathbf{s}$-Eulerian polynomials have only real roots},
Forum Math.\ Sigma \textbf{3} (2015), e1, 1--14.

\bibitem{Simion1984}
R.~Simion,
\emph{A multiindexed Sturm sequence of polynomials and unimodality of certain combinatorial sequences},
J.\ Combin.\ Theory Ser.\ A \textbf{36} (1984), 15--22.

\bibitem{TaNgo2026}
T.~T.~H. An and N.~Q. Hoan,
\emph{Multiple Eulerian polynomials},
J.\ Anal.\ \textbf{34} (2026), 21--34.

\end{thebibliography}

\end{document}
